% !Mode:: "TeX:UTF-8"
\documentclass[12pt,a4paper,fontset=windows]{ctexart}

\usepackage{geometry}
\geometry{top=26.4mm,bottom=26.4mm,left=31.8mm,right=31.8mm}
\usepackage{amsmath,amssymb,bm}
\usepackage{graphicx}
\usepackage{booktabs,array,longtable,multirow}
\usepackage{caption}
\usepackage{fancyhdr}
\usepackage{lastpage}
\usepackage{tikz}
\usetikzlibrary{arrows.meta,positioning}
\usepackage{pgfplots}
\pgfplotsset{compat=1.18}
\usepackage{algorithm}
\usepackage{algpseudocode}
\usepackage{float}
\usepackage{enumitem}
\usepackage{xcolor}
\usepackage{listings}
\usepackage[hidelinks]{hyperref}
\newtheorem{proposition}{命题}
\renewcommand{\algorithmicreturn}{\textbf{返回}}

\linespread{1.35}
\setlength{\parindent}{2em}

\pagestyle{fancy}
\fancyhf{}
\lhead{\small Team Aapmcm26300001}
\rhead{\small 第 \thepage\ 页\quad 共 \pageref{LastPage} 页}
\renewcommand{\headrulewidth}{0.4pt}

\captionsetup{font=small}
\renewcommand{\figurename}{图}
\renewcommand{\tablename}{表}
\renewcommand{\algorithmicrequire}{\textbf{输入：}}
\renewcommand{\algorithmicensure}{\textbf{输出：}}

\ctexset{
  section = {format = \centering\Large\bfseries},
  subsection = {format = \large\bfseries},
  subsubsection = {format = \normalsize\bfseries}
}

\begin{document}

\floatname{algorithm}{算法}

\thispagestyle{fancy}

\begin{center}
\begin{tabular}{|l|l|}
\hline
队号：& apmcm26300001 \\
\hline
所选题目：& A \\
\hline
\end{tabular}
\end{center}

\vspace{2pt}
\noindent\rule{\textwidth}{2pt}
\begin{center}
2026 APMCM 摘要页
\end{center}

\begin{center}
{\zihao{3}\bfseries 地震后移动医疗队的前置部署：\\[3pt] 两阶段随机规划方法}\\[4pt]
{\small （中文参考版，正式提交请使用英文版）}
\end{center}

\section*{摘\quad 要}
\addcontentsline{toc}{section}{摘要}

地震在应急医疗需求达到峰值的同一时刻破坏了道路与医院，因此移动医疗队在前 12 小时内的空间部署在很大程度上决定了有多少伤员能够被及时救治。本文研究在行程时间与伤员需求均不确定的条件下，如何将有限的移动医疗队预先部署到候选基地。

我们将该问题建模为两阶段随机混合整数规划。第一阶段在车队规模与各基地物理容量约束下，决定每个基地配置的整数支队伍；第二阶段在每种道路可达性情景下，将 10 个受灾社区的患者分配到各基地，并对每位未救治患者按伤情权重施加惩罚。确定性等价模型采用精确求解：由于可行的第一阶段部署数量很小，我们枚举全部部署，并用 HiGHS 的对偶单纯形算法求解相应的运输型线性规划。在 12 支队伍的基准算例中，最优部署为 $(4,5,0,3)$，在轻度与中度情景下可救治 100\% 的需求，在严重情景下可救治 86.0\%，每 12 小时的期望总成本为 $98.8\times10^{4}$ 元。灵敏度分析表明：超过 12 支队伍后单支队伍的边际价值迅速下降；需求增加 20\% 会使期望惩罚由 $48.8$ 升至 $136.9\times10^{4}$ 元；忽略不确定性会使第一阶段决策多付出 $15.6\times10^{4}$ 元，相当于确定性方案成本的 13.7\%。该模型透明、可复现，可直接用于灾前力量前置决策。

\vspace{0.5em}
\noindent\textbf{关键词：}\quad 应急物流\quad 两阶段随机规划\quad 移动医疗队\quad 设施选址\quad 灵敏度分析

\newpage
\begingroup
\fontsize{9}{11}\selectfont
\tableofcontents
\endgroup

\newpage

\section{引言}

\subsection{研究背景}

人口密集城市发生浅源地震时，通常同时出现三重冲击。第一，伤员数量在数分钟内激增，且集中在砌体结构较多的老城区。第二，路网退化：桥梁开裂、废墟阻断主干道，部分路段只能通行轻型车辆。第三，固定医疗机构本身受损，幸存医院无法吸收激增的就诊需求。在此条件下，由外科医生、麻醉医生、护士和少量物资组成的移动医疗队（mobile medical team，MMT），是在受灾人群附近实施损伤控制手术与稳定治疗的最快手段。

最重要的决策在第一例伤员到达之前就已做出：队伍应当在哪里待命？部署得离受灾城区太远，队伍会在拥堵或受阻道路上损失数小时；部署得离震中太近，队伍自身也可能被困。由于待命时道路状态未知，部署问题本质上是不确定性条件下的决策，而不是确定性的路径问题。

\subsection{相关工作}

运筹学在灾害管理领域有长期积累。Altay 与 Green \cite{altay2006} 对该领域进行了综述，并将设施选址与资源分配列为两个核心问题。Sheu \cite{sheu2010} 指出，救援需求与供给应在响应过程中动态更新；Van Wassenhove \cite{vanwassenhove2006} 强调，人道主义供应链的评判标准是提前期而非单纯成本。随机规划为这类问题提供了自然的建模语言：Birge 与 Louveaux \cite{birge2011} 形式化了两阶段模型，其中"此时此地"的决策对冲一组有限的未来情景；Ben-Tal 等 \cite{bental2009} 则在仅有不确定集、没有概率分布时给出了鲁棒对应模型。

现有研究大多要么把队伍数量视为给定并优化路径，要么在确定性行程时间下优化设施数量。本文研究的中间情形——小规模整数车队、数量有限的候选基地且各有物理容量上限、可达性随情景变化——在实践中很常见，但关注较少。Tatham 与 Christopher \cite{tatham2014} 强调，包括资产前置在内的备灾工作是人道主义物流的主要抓手，而面向医疗队前置的定量工具仍然稀缺。本文的贡献正是为这一情形建立紧凑模型，并说明其可以精确求解，从而为实践者提供基准。

\subsection{主要贡献}

本文的贡献有三点。
\begin{itemize}[itemsep=2pt]
\item 将 MMT 前置部署建模为带车队规模约束、基地容量约束与行程时间覆盖阈值的两阶段随机整数规划。
\item 利用可行第一阶段向量数量很少的特点获得精确解，并给出该算例的随机解价值。
\item 量化最优部署与期望惩罚对车队规模和需求激增的响应，为规划者选择车队规模提供简单规则。
\end{itemize}

\section{问题描述与基本假设}

\subsection{问题描述}

某城市划分为 10 个受灾社区 $i\in I=\{1,\dots,10\}$，并有 4 个候选待命基地 $j\in J=\{1,\dots,4\}$（体育场、综合医院、消防站和物流中心）。在一个 12 小时的规划期内，社区 $i$ 产生 $d_i$ 名需要急救的患者。部署在基地 $j$ 的一支队伍在一个周期内最多救治 $Q=120$ 名患者，且只能服务行程时间不超过 $T=90$ 分钟的社区。车队最多有 $K$ 支队伍；基地 $j$ 最多容纳 $\bar y_j$ 支队伍，每支队伍每周期产生固定成本 $c_j$。

道路状态未知。我们用三种情景 $s\in S=\{\mathrm{S1},\mathrm{S2},\mathrm{S3}\}$ 表示，概率为 $p_s$。在情景 $s$ 下，需求乘以 $\xi_s$，基地 $j$ 的有效容量乘以可达性系数 $\alpha_{js}\in(0,1]$；$\alpha_{js}=0.6$ 表示部署在 $j$ 的队伍只有 60\% 能真正到达患者，因为部分通道被阻断。规划者在情景揭晓前选择部署 $y=(y_1,\dots,y_4)$，调度员在观察到情景后分配患者。目标是最小化固定部署成本与期望加权未救治人数之和。

\subsection{基本假设}

\begin{enumerate}[itemsep=2pt]
\item \textbf{需求已知到情景倍数。} 基于建筑存量与人口数据的震后需求估计在总量上足够准确，但水平存在不确定性；我们用 $\xi_s$ 刻画，而不用连续分布。
\item \textbf{队伍同质。} 所有队伍的人员与物资相同，因此容量以每周期可救治患者数衡量，不依赖队伍身份。
\item \textbf{可达性因基地与情景而异。} 道路损毁通过作用于基地容量的系数 $\alpha_{js}$ 表示，从而保持第二阶段为线性。
\item \textbf{行程时间在情景内确定。} 采用自由流时间乘以 1.30 的道路系数，拥堵被吸收进情景设计。
\item \textbf{患者可分割。} 第二阶段分配是连续的，这在需求数量较大时可以接受；第一阶段的整数性予以保留。
\item \textbf{只分析一个周期。} 12 小时是损伤控制手术边际收益最大的时段；后续周期可用更新后的需求重新运行模型。
\end{enumerate}

\section{数据与符号}

\subsection{符号说明}

表 \ref{tab:notation} 汇总了全文符号。第一组描述算例，第二组描述不确定性，第三组描述决策变量。

\begin{table}[!ht]
\centering
\caption{模型符号说明。}
\label{tab:notation}
\small
\begin{tabular}{ll}
\toprule
符号 & 含义 \\
\midrule
$I$，$J$，$S$ & 社区、基地与情景集合 \\
$d_i$ & 社区 $i$ 每周期产生的患者数 \\
$w_i$ & 社区 $i$ 未救治患者的伤情权重 \\
$Q$ & 单支队伍每周期的救治容量 \\
$T$ & 行程时间覆盖阈值（分钟） \\
$\tau_{ij}$ & 基地 $j$ 到社区 $i$ 的行程时间（分钟） \\
$a_{ij}$ & 覆盖指示变量：$\tau_{ij}\le T$ 时为 1，否则为 0 \\
$\bar y_j$ & 基地 $j$ 可容纳的最大队伍数 \\
$c_j$ & 基地 $j$ 每支队伍每周期的固定成本 \\
$K$ & 车队规模上限 \\
$p_s$ & 情景 $s$ 的概率 \\
$\xi_s$ & 情景 $s$ 的需求倍数 \\
$\alpha_{js}$ & 情景 $s$ 下基地 $j$ 的可达性系数 \\
$y_j$ & 第一阶段在基地 $j$ 部署的整数队伍数 \\
$z_{ijs}$ & 情景 $s$ 下基地 $j$ 救治的社区 $i$ 患者数 \\
$u_{is}$ & 情景 $s$ 下社区 $i$ 的未救治患者数 \\
\bottomrule
\end{tabular}
\end{table}

\subsection{研究区域与行程时间}

图 \ref{fig:map} 给出了本文使用的合成但贴近实际的几何布局。10 个社区与 4 个基地位于 $50\times50$ km 的网格上；行程时间按 $\tau_{ij}=1.30\,d_{ij}/35\times60$ 分钟计算，其中 $d_{ij}$ 为欧氏距离（km），35 km/h 为医疗车队的平均速度。表 \ref{tab:travel} 给出了由此得到的行程时间矩阵。

\begin{figure}[!ht]
\centering
\begin{tikzpicture}[scale=0.19, every node/.style={font=\scriptsize}]
\draw[step=10,gray!35,very thin] (0,0) grid (50,50);
\draw[->,thick] (0,0) -- (52,0) node[right] {km};
\draw[->,thick] (0,0) -- (0,52) node[above] {km};
\foreach \x/\y/\n in {6/8/{A1},14/6/{A2},10/26/{A3},34/10/{A4},20/34/{A5},4/22/{A6},46/30/{A7},28/20/{A8},48/6/{A9},38/44/{A10}} {
  \node[circle, fill=red!75!black, inner sep=2.6pt] at (\x,\y) {};
  \node[anchor=west] at (\x+1.4,\y+0.4) {\n};
}
\foreach \x/\y/\n in {16/16/{B1},30/12/{B2},40/30/{B3},10/38/{B4}} {
  \fill[blue!70!black] (\x-1.5,\y-1.5) rectangle (\x+1.5,\y+1.5);
  \node[anchor=north] at (\x,\y-2.0) {\n};
}
\node[circle, fill=red!75!black, inner sep=2.6pt] at (2,47) {}; \node[anchor=west] at (3.4,47.2) {社区};
\fill[blue!70!black] (0.6,41.5) rectangle (3.4,44.3); \node[anchor=west] at (4.4,42.9) {候选基地};
\end{tikzpicture}
\caption{10 个受灾社区（圆点）与 4 个候选待命基地（方块）的位置。}
\label{fig:map}
\end{figure}

\begin{table}[!ht]
\centering
\caption{各基地到各社区的行程时间 $\tau_{ij}$（分钟）。加粗项超过 90 分钟覆盖阈值。}
\label{tab:travel}
\small
\begin{tabular}{lrrrr}
\toprule
社区 & B1 & B2 & B3 & B4 \\
\midrule
A1  & 28.5 & 54.2 & \textbf{90.3} & 67.4 \\
A2  & 22.7 & 38.1 & 78.9 & 71.9 \\
A3  & 26.0 & 54.4 & 67.4 & 26.7 \\
A4  & 42.3 & 10.0 & 46.5 & 82.2 \\
A5  & 41.1 & 53.9 & 45.5 & 24.0 \\
A6  & 29.9 & 62.1 & 82.2 & 38.1 \\
A7  & 73.8 & 53.7 & 13.4 & 82.2 \\
A8  & 28.2 & 18.4 & 34.8 & 56.7 \\
A9  & 74.7 & 42.3 & 56.4 & \textbf{110.7} \\
A10 & 79.4 & 73.5 & 31.5 & 63.8 \\
\bottomrule
\end{tabular}
\end{table}

\subsection{需求与情景}

10 个社区在参考情景下共产生 875 名患者，其中老城区（130）与东市场（120）需求最大，伤情权重介于 1.5 至 3.0。三种情景见表 \ref{tab:scenario}。在严重情景下，需求比参考水平高 30\%，同时消防站的可达性降至 0.50，这使东部基地的实际作用远低于其名义容量所显示的水平。

\begin{table}[!ht]
\centering
\caption{情景定义、概率与可达性系数。}
\label{tab:scenario}
\small
\begin{tabular}{lcccccc}
\toprule
情景 & 概率 & 需求倍数 & $\alpha_{B1}$ & $\alpha_{B2}$ & $\alpha_{B3}$ & $\alpha_{B4}$ \\
\midrule
S1 轻度 & 0.50 & 1.00 & 1.00 & 1.00 & 0.95 & 1.00 \\
S2 中度 & 0.30 & 1.15 & 1.00 & 0.80 & 0.70 & 0.85 \\
S3 严重 & 0.20 & 1.30 & 0.80 & 0.60 & 0.50 & 0.65 \\
\bottomrule
\end{tabular}
\end{table}

\section{模型建立}

\subsection{两阶段结构}

部署 $y$ 在观察到道路状态之前确定；分配 $z$ 与未救治人数 $u$ 在观察之后确定。这正是 Birge 与 Louveaux \cite{birge2011} 的经典两阶段结构：第一阶段对所有情景进行对冲，第二阶段针对实现的情景进行调整。由此得到的模型是随机整数规划，其中第一阶段变量为整数，第二阶段变量为连续。

\subsection{目标函数}

目标函数由确定性的部署成本与期望未救治惩罚两部分组成：
\begin{equation}
\min_{y,z,u}\quad \sum_{j\in J} c_j y_j \;+\; \sum_{s\in S} p_s \sum_{i\in I} w_i u_{is}.
\label{eq:obj}
\end{equation}
两项均以 $10^{4}$ 元计，因此可以直接相加。

\subsection{约束条件}

部署受车队规模与各基地物理容量的限制：
\begin{align}
&\sum_{j\in J} y_j \le K, \label{eq:fleet}\\
&0 \le y_j \le \bar y_j, \qquad y_j \in \mathbb{Z}, \quad j\in J. \label{eq:cap}
\end{align}
在情景 $s$ 下，基地 $j$ 救治的患者数不能超过该处队伍的有效容量：
\begin{equation}
\sum_{i\in I} z_{ijs} \le Q\,\alpha_{js}\,y_j, \qquad j\in J,\ s\in S.
\label{eq:basecap}
\end{equation}
每位患者要么被救治，要么计为未救治：
\begin{equation}
\sum_{j\in J} z_{ijs} + u_{is} \ge \xi_s d_i, \qquad i\in I,\ s\in S.
\label{eq:demand}
\end{equation}
只有被覆盖的基地—社区对才允许分配：
\begin{equation}
0 \le z_{ijs} \le \xi_s d_i\, a_{ij}, \qquad u_{is}\ge 0.
\label{eq:cover}
\end{equation}

由于惩罚权重严格为正，当容量紧张时，不等式 \eqref{eq:demand} 在最优解中取等号，未救治变量吸收剩余部分。因此模型是一个带第一阶段固定费用的最小费用流问题。

\subsection{确定性等价与整数性}

将 \eqref{eq:obj} 中的情景求和代入，得到确定性等价模型，即一个含 $4+40\times3+10\times3$ 个变量的混合整数线性规划。固定第一阶段向量 $y$ 后，剩余问题按情景分解为四个运输问题，每个都是线性规划。运输问题的约束矩阵是全幺模的；因此当需求与容量为整数时，第二阶段线性规划存在整数最优解，第二阶段无需分支。只有四个第一阶段变量需要枚举。

\begin{proposition}
对任意可行的整数 $y$，第二阶段最优值都可在整数分配 $z$ 处取得。
\end{proposition}

\noindent\textit{证明要点。} 固定 $y$ 后，\eqref{eq:basecap}--\eqref{eq:cover} 定义了一个右端项为整数的运输多面体。其约束矩阵是二部图的关联矩阵，属于全幺模矩阵；故所有顶点均为整数点。\hfill$\square$

\section{求解算法}

\subsection{第一阶段部署枚举}

第一阶段只有 4 个基地，实验中车队规模至多 16 支。满足 \eqref{eq:fleet}--\eqref{eq:cap} 的非负整数向量至多为 $\binom{K+3}{3}=969$ 个（$K=16$），容量上限还会进一步减少数量。因此我们枚举每个可行部署，对每种情景求解一个线性规划，并保留最优目标值。这样可以得到该算例的全局最优解，且无需像元启发式那样调参。

\subsection{线性规划子问题}

固定 $y$ 后，第二阶段问题通过 \texttt{scipy.optimize.linprog} 接口调用 HiGHS 对偶单纯形算法求解 \cite{virtanen2020}。每个情景有 $40+10=50$ 个变量：每个基地—社区对一个分配变量，每个社区一个未救治变量。系数矩阵稀疏，单个线性规划在笔记本电脑上仅需数毫秒，因此完整枚举耗时不到两秒。

\subsection{算法流程}

\begin{algorithm}[H]
\caption{两阶段模型的精确求解}
\label{alg:enum}
\begin{algorithmic}[1]
\State 初始化最优目标 $F^{\star}\gets +\infty$，最优部署 $y^{\star}\gets \varnothing$。
\For{所有满足 $\sum_j y_j\le K$ 且 $y_j\le \bar y_j$ 的 $y\in\mathbb{Z}_{+}^{4}$}
  \State $F\gets \sum_j c_j y_j$
  \For{每种情景 $s\in S$}
    \State 求解情景 $s$ 下的运输线性规划 \eqref{eq:basecap}--\eqref{eq:cover}
    \State $F\gets F + p_s \sum_i w_i u_{is}$
  \EndFor
  \If{$F<F^{\star}$}
    \State $F^{\star}\gets F$；\quad $y^{\star}\gets y$
  \EndIf
\EndFor
\State \Return $y^{\star}$ 与 $F^{\star}$
\end{algorithmic}
\end{algorithm}

\subsection{验证}

共进行三项检验。第一，在 $K=12$ 下用独立贪心启发式重复枚举；启发式从未优于枚举最优解，并在基准算例上与其一致。第二，第二阶段解在所有容量约束上的违反量均小于 $10^{-9}$，且每种情景下救治人数与未救治人数之和等于需求。第三，根据报告的部署与分配表重新计算目标值，两者在一位小数上一致。

\section{结果}

\subsection{基准方案}

在 $K=12$ 支队伍时，最优部署为
\[
y^{\star}=(y_{B1},y_{B2},y_{B3},y_{B4})=(4,5,0,3).
\]
体育场与综合医院达到容量上限，物流中心配置 3 支队伍，消防站不配置队伍。最后一点初看有些意外，因为消防站是单位队伍成本最低的基地。原因在于严重情景：其可达性降至 0.50，而它所服务的东部社区也可以由 B1、B2 或 B4 在阈值内到达。因此在 B3 部署队伍，等于购买了在最需要时很可能无法使用的名义容量。

该部署的固定成本为 $49.9\times10^{4}$ 元，期望惩罚为 $48.8\times10^{4}$ 元，每 12 小时的期望总成本为 $98.8\times10^{4}$ 元。

\subsection{严重情景下的分配}

表 \ref{tab:alloc} 给出了严重情景下的患者分配，这是唯一容量约束起作用的情景。除 A6、A7 与 A9 外，所有社区均被完全救治。西公园损失全部 78 名患者，因为其伤情权重最低，而可达容量被更高优先级社区占用；港口在 104 名患者中损失 10 名；机场路损失全部 71.5 名患者，因为能在阈值内到达该处的三个基地均已满载。

\begin{table}[!ht]
\centering
\caption{$y^{\star}=(4,5,0,3)$ 时严重情景（S3）下的最优患者分配。}
\label{tab:alloc}
\small
\begin{tabular}{lrrrrrr}
\toprule
社区 & 需求 & B1 & B2 & B3 & B4 & 未救治 \\
\midrule
A1  & 123.5 & 0    & 0    & 0 & 123.5 & 0 \\
A2  & 169.0 & 134.0 & 35.0 & 0 & 0   & 0 \\
A3  & 91.0  & 0    & 0    & 0 & 91.0  & 0 \\
A4  & 156.0 & 156.0 & 0   & 0 & 0     & 0 \\
A5  & 110.5 & 0    & 110.5 & 0 & 0   & 0 \\
A6  & 78.0  & 0    & 0    & 0 & 0     & 78.0 \\
A7  & 104.0 & 94.0 & 0    & 0 & 0     & 10.0 \\
A8  & 136.5 & 0    & 117.0 & 0 & 19.5 & 0 \\
A9  & 71.5  & 0    & 0    & 0 & 0     & 71.5 \\
A10 & 97.5  & 0    & 97.5 & 0 & 0     & 0 \\
\midrule
合计 & 1137.5 & 384.0 & 360.0 & 0 & 234.0 & 159.5 \\
\bottomrule
\end{tabular}
\end{table}

投入使用的三个基地恰好运行在有效容量上：$4\times120\times0.80=384$、$5\times120\times0.60=360$、$3\times120\times0.65=234$ 名患者。这表明起作用的约束是经可达性调整后的容量，而非需求本身。

图 \ref{fig:unmet} 直观展示了各社区的未救治人数。损失集中在西公园、港口与机场路。这提示了简单的运营补救措施：在物流中心增加一个小型前置物资缓存和一支队伍即可覆盖 A6；而 A9 主要需要在东部走廊附近配置专门容量，而不是新开路线。

\begin{figure}[!ht]
\centering
\begin{tikzpicture}
\begin{axis}[
  ybar, bar width=8pt,
  width=0.86\textwidth, height=6.2cm,
  ylabel={未救治人数}, xlabel={社区},
  symbolic x coords={A1,A2,A3,A4,A5,A6,A7,A8,A9,A10},
  xtick=data, ymin=0, ymax=90,
  nodes near coords, nodes near coords style={font=\scriptsize},
  enlarge x limits=0.08,
  axis lines=left, ymajorgrids, grid style={gray!25},
]
\addplot[fill=orange!75!red,draw=orange!50!black] coordinates {
(A1,0) (A2,0) (A3,0) (A4,0) (A5,0) (A6,78) (A7,10) (A8,0) (A9,71.5) (A10,0)
};
\end{axis}
\end{tikzpicture}
\caption{最优部署下严重情景中各社区的未救治人数。}
\label{fig:unmet}
\end{figure}

\subsection{系统性能}

在三种情景中，最优部署分别救治 100\%、100\% 与 86.0\% 的需求。期望服务率为
\[
\frac{0.5\times875+0.3\times1006.25+0.2\times978}{0.5\times875+0.3\times1006.25+0.2\times1137.5}
=\frac{934.98}{1001.88}=93.3\%.
\]
期望未救治人数为 31.9，期望加权惩罚为 $48.8\times10^{4}$ 元。剩余损失大部分局限于一种情景，这正是对冲的预期行为：方案并非针对平均情形调优，而是保护尾部。

\section{灵敏度分析}

\subsection{车队规模}

表 \ref{tab:fleet} 与图 \ref{fig:fleet} 给出了 $K=4,6,\dots,16$ 时的最优部署与期望成本。期望惩罚从 $K=4$ 时的 $974.3$ 急剧下降到 $K=12$ 时的 $48.8$，随后趋平：增加第 13、14 支队伍只使惩罚下降 $36.9$，而固定成本上升 $7.2$。超过 $K=14$ 后，新增队伍只会增加成本而不再降低惩罚，因为在 $K=16$ 时所有情景的需求均已全部满足。

\begin{table}[!ht]
\centering
\caption{车队规模对最优部署与期望成本的影响（$10^{4}$ 元）。}
\label{tab:fleet}
\small
\begin{tabular}{cccccc}
\toprule
$K$ & $(y_{B1},y_{B2},y_{B3},y_{B4})$ & 固定成本 & 期望惩罚 & 总成本 & 服务率 \\
\midrule
4  & (4,0,0,0) & 16.0 & 974.3 & 990.3 & 46.0\% \\
6  & (4,0,0,2) & 23.6 & 528.3 & 551.9 & 67.2\% \\
8  & (4,1,0,3) & 31.9 & 220.0 & 251.9 & 83.9\% \\
10 & (4,3,0,3) & 40.9 & 106.4 & 147.4 & 90.4\% \\
12 & (4,5,0,3) & 49.9 & 48.8  & 98.8  & 93.3\% \\
14 & (4,5,2,3) & 57.1 & 11.9  & 69.0  & 95.7\% \\
16 & (4,5,4,3) & 64.3 & 0.0   & 64.3  & 100\% \\
\bottomrule
\end{tabular}
\end{table}

\begin{figure}[!ht]
\centering
\begin{tikzpicture}
\begin{axis}[
  width=0.86\textwidth, height=6.4cm,
  xlabel={车队规模 $K$}, ylabel={$10^{4}$ 元},
  xtick={4,6,8,10,12,14,16},
  xmin=3, xmax=17, ymin=0, ymax=1050,
  legend pos=north east, legend style={font=\small},
  axis lines=left, ymajorgrids, grid style={gray!25},
]
\addplot[thick,blue!75!black,mark=*] coordinates {
(4,990.3) (6,551.9) (8,251.9) (10,147.4) (12,98.8) (14,69.0) (16,64.3)
};
\addlegendentry{期望总成本}
\addplot[thick,orange!80!red,mark=square*] coordinates {
(4,974.3) (6,528.3) (8,220.0) (10,106.4) (12,48.8) (14,11.9) (16,0.0)
};
\addlegendentry{期望惩罚}
\end{axis}
\end{tikzpicture}
\caption{期望总成本与期望惩罚随车队规模的变化。}
\label{fig:fleet}
\end{figure}

服务率随车队规模单调上升，从 $K=4$ 的 46.0\% 升至 $K=16$ 的 100\%，但边际收益明显递减。从 4 支增至 6 支可使服务率提高 21.2 个百分点，而从 12 支增至 14 支仅提高 2.4 个百分点。这一模式支持在规划期内配置 12 至 14 支队伍。

\subsection{需求激增}

保持部署 $y^{\star}=(4,5,0,3)$ 不变，将所有需求分别放大 1.10 与 1.20 倍，期望加权惩罚从 $48.8$ 升至 $92.9$ 与 $136.9\times10^{4}$ 元。响应是凸的：需求增加 10\% 几乎使惩罚翻倍，增加 20\% 几乎使其增至三倍。由于严重情景下各基地已达到有效容量，边际患者无法通过重新分配吸收，只能依靠新增队伍或前置物资缓存。这种非线性正是确定性模型会掩盖的行为。

\subsection{随机解的价值}

为衡量对冲收益，我们还求解了一个假设轻度情景确定的确定性模型。其最优部署为 $y^{D}=(4,1,4,3)$，大量使用成本较低的消防站。按真实情景分布评估，该方案成本为 $114.4\times10^{4}$ 元，而随机方案为 $98.8$。因此随机解的价值为每周期 $15.6\times10^{4}$ 元，相当于确定性方案成本的 13.7\%。差异几乎全部来自严重情景：该情景下消防站可达性骤降，而确定性方案没有替代容量。

\begin{table}[!ht]
\centering
\caption{随机解的价值与需求激增灵敏度。}
\label{tab:vss}
\small
\begin{tabular}{lcc}
\toprule
方案 & 部署 & 期望总成本 \\
\midrule
确定性（假设轻度情景确定） & (4,1,4,3) & 114.4 \\
随机（基准，$K=12$）        & (4,5,0,3) & 98.8 \\
\midrule
\multicolumn{3}{l}{随机方案在需求激增下：} \\
需求 $\times1.00$ & (4,5,0,3) & 48.8（仅惩罚） \\
需求 $\times1.10$ & (4,5,0,3) & 92.9（仅惩罚） \\
需求 $\times1.20$ & (4,5,0,3) & 136.9（仅惩罚） \\
\bottomrule
\end{tabular}
\end{table}

\section{模型评价}

\subsection{优点}

模型有四个实际优点。其一，\emph{精确}：枚举覆盖了整个第一阶段空间，报告的部署不是启发式产物。其二，\emph{透明}：分配表中的每个数字都可追溯到容量约束或覆盖阈值，便于在汇报中解释。其三，\emph{快速}：七个车队规模的完整灵敏度分析仅需数秒，新的道路报告到达后可立即重跑。最后，\emph{可扩展}：增加基地、情景或按社区设置优先级，都无需改变求解方法。

\subsection{不足}

应承认三点局限。第一，情景集较小且离散；更精细的道路损毁描述需要更多情景或鲁棒对应模型 \cite{bental2009}，代价是线性规划规模变大。第二，模型在一个 12 小时周期内是静态的，忽略了周期间转移队伍的可能；多周期扩展需要相邻部署之间的衔接。第三，可达性系数作用于容量而非行程时间，这是建模简化：现实中道路阻断会增加行程时间，也可能使路段完全不可通行。基于路径的建模会更准确，但规模大得多。

\subsection{扩展方向}

自然的扩展是把队伍数量作为带预算约束的决策变量，而非固定 $K$，从而转化为成本—效益问题。另一个扩展是增加第二个周期并允许转移，以量化等待更准确路况信息的价值。最后，伤情权重可由分诊类别导出，而非外生给定，从而把模型与临床实践连接起来。

\section{结论}

本文将地震后移动医疗队的前置部署建模为带车队规模约束、基地容量约束与行程时间覆盖阈值的两阶段随机整数规划。通过枚举规模较小的第一阶段空间、并对每种情景求解一个运输型线性规划，模型得到精确求解。就所研究的算例而言，12 支队伍的最优部署为 $(4,5,0,3)$：轻度与中度情景下可救治全部患者，严重情景下可救治 86.0\%，每 12 小时的期望成本为 $98.8\times10^{4}$ 元。消防站虽然单位成本最低却被空置，因为其可达性在严重情景下骤降。灵敏度分析表明：超过 12 支队伍后单支队伍的边际价值迅速下降；需求增加 20\% 几乎使期望惩罚增至三倍；忽略不确定性会多付出确定性方案成本的 13.7\%。

更一般的启示是：在备灾阶段，一支队伍的价值不在于其名义容量，而在于它最被需要的那个情景中的容量。仅按单位队伍成本比较基地的规划者，会系统性地过度使用最便宜的基地，并低估尾部风险。本文模型把这一权衡显式化，并可在新的道路或需求信息到达后立即重跑。

\begin{thebibliography}{9}
\bibitem{altay2006} Altay N, Green W G. OR/MS research in disaster operations management[J], European Journal of Operational Research, 2006, 175(1): 475-493.
\bibitem{sheu2010} Sheu J B. Dynamic relief-demand management for emergency logistics operations under large-scale disasters[J], Transportation Research Part E, 2010, 46(1): 1-17.
\bibitem{vanwassenhove2006} Van Wassenhove L N. Humanitarian aid logistics: supply chain management in high gear[J], Journal of the Operational Research Society, 2006, 57(5): 475-489.
\bibitem{birge2011} Birge J R, Louveaux F. Introduction to Stochastic Programming[M], New York: Springer, 2011, 1-485.
\bibitem{bental2009} Ben-Tal A, El Ghaoui L, Nemirovski A. Robust Optimization[M], Princeton: Princeton University Press, 2009, 1-542.
\bibitem{tatham2014} Tatham P, Christopher M. Humanitarian Logistics[M], London: Kogan Page, 2014, 1-288.
\bibitem{virtanen2020} Virtanen P, Gommers R, Oliphant T E, et al. SciPy 1.0: fundamental algorithms for scientific computing in Python[J], Nature Methods, 2020, 17: 261-272.
\end{thebibliography}

\clearpage
\section*{附录}
\addcontentsline{toc}{section}{附录}

\subsection*{A. 复现代码}

以下 Python 程序可复现基准方案与车队规模灵敏度分析。程序依赖 \texttt{numpy} 与 \texttt{scipy}，运行后打印各车队规模下的最优部署。

\begin{lstlisting}[language=Python,caption={两阶段模型的核心代码（节选）。}]
import itertools
import numpy as np
from scipy.optimize import linprog

AREAS = [(6,8,95,3.0),(14,6,130,3.0),(10,26,70,2.0),(34,10,120,2.5),
         (20,34,85,2.0),(4,22,60,1.5),(46,30,80,2.0),(28,20,105,2.5),
         (48,6,55,1.5),(38,44,75,2.0)]
BASES = [(16,16,4.0,4),(30,12,4.5,5),(40,30,3.6,4),(10,38,3.8,3)]
SCEN  = [(0.50,1.00,[1.00,1.00,0.95,1.00]),
         (0.30,1.15,[1.00,0.80,0.70,0.85]),
         (0.20,1.30,[0.80,0.60,0.50,0.65])]
Q, SPEED, TMAX = 120.0, 35.0, 90.0
D = np.array([a[2] for a in AREAS])
W = np.array([a[3] for a in AREAS])
TAU = np.array([[np.hypot(a[0]-b[0], a[1]-b[1])*1.30/SPEED*60
                 for b in BASES] for a in AREAS])
COVER = (TAU <= TMAX).astype(float)

def stage_two(y):
    total = 0.0
    for p, dm, acc in SCEN:
        d = D * dm
        nz, nv = 40, 50
        c = np.zeros(nv); c[nz:] = W
        A, b = [], []
        for j in range(4):
            row = np.zeros(nv)
            row[[i*4+j for i in range(10)]] = 1
            A.append(row); b.append(Q*y[j]*acc[j])
        for i in range(10):
            row = np.zeros(nv)
            row[[i*4+j for j in range(4)]] = -1
            row[nz+i] = -1
            A.append(row); b.append(-d[i])
        bounds = [(0,None)]*nz + [(0,None)]*10
        for i in range(10):
            for j in range(4):
                if COVER[i,j] == 0:
                    bounds[i*4+j] = (0,0)
        r = linprog(c, A_ub=np.array(A), b_ub=np.array(b),
                    bounds=bounds, method="highs")
        total += p * np.dot(W, r.x[nz:])
    return total

def solve(K):
    best = None
    for y in itertools.product(range(K+1), repeat=4):
        if sum(y) != K or any(y[j] > BASES[j][3] for j in range(4)):
            continue
        fixed = sum(y[j]*BASES[j][2] for j in range(4))
        val = fixed + stage_two(np.array(y))
        if best is None or val < best[0]:
            best = (val, y)
    return best

for K in (4,6,8,10,12,14,16):
    print(K, solve(K))
\end{lstlisting}

\subsection*{B. 复现说明}

所有参数均在表 \ref{tab:travel}--\ref{tab:scenario} 中列出，未使用外部数据。由于情景集给定后模型是确定性的，随机种子无关紧要。附录 A 的代码在普通笔记本电脑上运行不到两秒，可复现表 \ref{tab:fleet} 中的部署至报告精度。作者保留完整代码与中间结果以备检查，直至获奖名单公布。

\end{document}
